\chapter{钱不是财富，钱是一套记账系统}
\chaptermeta{01}

\begin{ledetext}
序章那 98 块只是两行数字。可老板真的把面端上来了。这一章问的就是这件事：一行数字凭什么换来一碗面。
\end{ledetext}

\section{有一块石头沉在海底，它照样是钱}

太平洋上有个岛叫雅浦。岛上的钱是石头。

石灰岩凿的圆盘，中间开个洞，大的直径三米多，重四吨。雅浦不产石灰岩，得划独木舟去四百公里外的帕劳采，再运回来。翻过船，死过人。所以一块石头值多少，不光看大小，还看它是怎么运回来的。

石头太重，搬不动。于是付款的时候石头不动。一户人家把石头付给另一户，做法是当着一群人的面说一句"这块归你了"。石头还立在原地，青苔照长，主人换了。

1910 年前后，人类学家弗内斯在岛上碰到一户特别富的人家。全村没人见过他家的钱。

那块巨石在运回来的路上遇了风暴，沉了。

船员作证：石头确实凿出来了，尺寸成色都好，就是现在躺在海底。全村认了。这块谁也没见过、也捞不上来的石头，被这家人当财产传了几代，还照样拿出来付账。

\begin{pullquote}{}
钱不是东西。钱是一群人共同维护的一段记忆。
\end{pullquote}

1932 年，法国怕美国放弃金本位，要求把存在纽约联邦储备银行的美元换成黄金，黄金不用运回巴黎。

美联储照办了。工作人员走进金库，把相应数量的金条从一个抽屉挪到另一个抽屉，换上写着"法国"的标签。

金条一米没动。第二天报纸开始报道美国黄金外流，美元承压。

弗里德曼后来点破：纽约金库里换标签，和雅浦岛给沉在海底的石头换主人，是同一件事。区别只在于我们这边有抽屉。

\begin{egbox}{举个栗子}
假设一个村子，十七户人家，没有一张钞票，村口挂着一块黑板。

李建国帮邻居修了三天屋顶，黑板上记一笔："王有福欠李建国 3 个工。"王有福给第三家修拖拉机，再记一笔。月底村会计做一次轧差，中间那些环节全划掉，只剩下谁真欠谁多少。

这个村子有钱吗。没有一张纸，但它有价格、有债务、有清算、有支付。

黑板就是货币系统，钞票是黑板的便携版。把村子放大到 14 亿人，黑板换成机房，村会计换成商业银行和央行。

\end{egbox}

\section{教科书上那个起源故事是编的}

故事你一定听过。很久以前人们物物交换，你有一头牛想换斧头，可拿斧头的人只想要鞋，交易做不成。这叫\term{双重巧合难题}{double coincidence of wants}\index{055@双重巧合难题}。为了解决它，人类发明了货币。

出自亚当·斯密 1776 年的《国富论》，写得漂亮，被抄了两百五十年。

它是编的。

人类学家找了一百多年，没找到任何一个靠物物交换运转的社会。1985 年剑桥的汉弗莱写过一句被引到烂的话：从来没有人描述过一个纯粹的物物交换经济，更别说从里面长出货币。

田野里能看到的物物交换，基本发生在陌生人和敌人之间。彼此不信任，才必须一手交钱一手交货。

村子内部靠赊账。你今天帮我收麦子，下个月我杀猪给你留条腿。谁欠谁多少全在人情里，攒到某个时候一起算。

考古的方向也对得上。美索不达米亚的泥板早在公元前三千年就密密麻麻记着"某人欠神庙多少大麦""折合多少银子"，而最早的铸币出现在公元前七世纪的吕底亚。\hlmark{记账比铸币早了两千多年。}

\begin{mythbox}{常见误解}
\mvfrow{误解}{"人类先物物交换，后来嫌麻烦发明了货币，所以钱本质上是一种方便的商品。"}
\mvfrow{事实}{顺序反了。\textbf{先有债务和记账，后有铸币。}铸币解决的是跟陌生人、跟远方、跟军队做交易的问题，不是从村口的以物易物里自然长出来的。"双重巧合难题"是一个功能论证，它说明货币解决了什么，它不是历史记录。因为伞能挡雨就断定伞是雨发明的，逻辑是一样的。}
诚实地补一句：细节上学界还在吵。格雷伯的《债：第一个五千年》把这套说法推给了大众，其中一些史料解读经济史学界并不买账。但"信用早于铸币"这条主干，现在很少有人反对。

\end{mythbox}

\section{三个功能，一个老大}

\begin{flowlist}
  \flowitem{01}{交换媒介\enterm{medium of exchange}}{你拿它换东西。最直观，也最容易被替代。贝壳、盐、烟草、方便面都干过这活儿。}
  \flowitem{02}{记账单位\enterm{unit of account}}{价格、工资、合同、房贷、债券票息，都用它标价。它是这个世界的度量衡。}
  \flowitem{03}{价值储藏\enterm{store of value}}{今天不花，留到明天还有用。三个功能里最脆的一个，通胀会慢慢啃掉它。}
\end{flowlist}

这三个不是平级的。\textbf{记账单位才是老大。}

一国货币真正死掉的那一刻，不是大家不再拿它买东西，是商家不再拿它标价。

恶性通胀的国家里，这两件事是分开发生的。价签早改成了美元，结账还收本币，按当天黑市汇率折。这时候本币只剩一个找零工具的身份，度量衡那把椅子已经让出去了。1923 年的德国走得更远：人们干脆用香烟和实物直接换。

反过来，只要工资、合同和债务还用某种货币标价，它就有极强的粘性。你欠银行的 300 万房贷是用人民币写的，你就得想办法弄到人民币。

\begin{egbox}{举个栗子}
1945 年，拉德福德写了篇论文，讲他待了两年的德国战俘营。营里的钱是香烟。

一罐牛肉 8 支烟，洗一件衬衫 2 支。有人在营区两头跑价差，有人开"咖啡馆"赊账供应，用烟结算。

红十字包裹一到，香烟供给暴增，所有东西的烟价当天上涨。这是通胀。断供几周，大家把烟抽掉了，物价开始跌。这是通缩。

2010 年代美国不少监狱禁烟，囚犯的通货换成了方便面。

香烟和方便面成为货币，不是因为好抽好吃。是因为大家开始拿它标价。

\end{egbox}

\section{老陈那张一百块凭什么}

老陈在昆明开面馆，二十六年了。他收人民币，不收比特币，也不收你写的欠条。

问他为什么，他会说这不废话吗。

不废话。他敢收，是因为四件事同时成立。

\begin{flowlist}
  \flowitem{01}{国家拿它收税}{最硬的一根。老陈每季度要交增值税和附加，雇的三个人要交个税，全得用人民币。他就算特别喜欢黄金，也得先把黄金换成人民币再去交。几亿人被迫弄到手的东西，天然有需求。}
  \flowitem{02}{法偿性\enterm{legal tender}}{法律规定用人民币偿还境内债务，债权人不得拒收。中国人民银行公开整治过拒收现金，罚过商户。这条的真实含义是：老陈手里的钱，法院认。}
  \flowitem{03}{网络效应}{老陈收人民币，因为他房东收、员工收、菜市场收。货币跟语言一样，用的人越多越好用，也越难改。你可以自己发明一种语言，但你没法让整条街的菜贩子学。}
  \flowitem{04}{央行盯着它的购买力}{货币是有质量的产品，质量就是购买力。2026 年 8 月，美联储把联邦基金利率维持在 \textbf{3.50\%–3.75\%}，年内第五次按兵不动，理由是通胀还高。它宁可让经济难受，也要保住这张纸的信用。这份信誉攒了几十年，砸掉要不了几年。}
\end{flowlist}

\begin{mythbox}{常见误解}
\mvfrow{误解}{"钱有价值是因为国家有黄金储备。国库里的金子没了，钱就废了。"}
\mvfrow{事实}{1971 年 8 月 15 日美国关掉黄金窗口之后，\textbf{世界上没有任何一种主要货币能兑换黄金}。中国人民银行不承诺给你换金条，美联储也不承诺。黄金今天是各国央行的储备资产之一，是压舱石，不是发钞的抵押物。撑住老陈那张一百块的，是上面那四条：税收、法律、网络、信誉。}
\end{mythbox}

\section{金本位回不去了}

怀念金本位的人通常这么说：钞票能换金子，政府就没法乱来。

前半句对。后半句的代价很大。

\begin{flowlist}
  \flowitem{01}{货币供给由挖矿速度定，不由经济需要定}{1873 到 1896 年，美国物价跌了二十多年。农民卖粮的钱一年比一年少，欠银行的债还是那个数字。1896 年布莱恩喊"你们不能把人类钉在黄金的十字架上"，喊的就是这个。反过来，加州淘金和南非金矿又把物价推上去。把货币的松紧交给地质运气，这不叫纪律，叫抽签。}
  \flowitem{02}{危机沿着固定汇率传染}{一国货币被挤兑，央行只能加息紧缩来保住兑换承诺，哪怕国内已经是大萧条。经济史里有条刺眼的规律：越早脱离金本位的国家，复苏越早。英国 1931 年 9 月脱钩，比死守到最后的国家早好几年爬出来。}
\end{flowlist}

1944 年的布雷顿森林是金本位的最后一版：美元按 35 美元一盎司锚黄金，别的货币锚美元。撑了不到三十年。1971 年，美国的黄金已经远远盖不住海外的美元债权，尼克松关了窗口。1973 年，主要货币开始浮动。

从那天起全世界的钱都是\term{法定货币}{fiat money}\index{013@法定货币}。fiat 是拉丁文，意思是"就这么定了"。

\begin{databox}{2026 数据：黄金回来了，但不是以你想的方式}
欧洲央行 2026 年 6 月的《欧元的国际角色》报告：截至 2025 年底，\textbf{黄金已经超过美国国债，成为全球央行的第一大储备资产}。

\begin{barfig}
  \barrow{黄金}{27}{27\%}
  \barrow{美国国债}{22}{22\%}
  \barrow{其他美元资产}{20}{20\%}
  \barrow{欧元资产}{15}{15\%}
\end{barfig}

黄金占比一年之间从 20\% 升到 27\%，美债从 25\% 降到 22\%。但另一个数字要一起读：\textbf{美元计价资产合计仍占 42\%}，还是最大的单一货币阵营。这是渐进的多元化，不是改朝换代。

其他口径：全球央行黄金储备超 \textbf{3.6 万吨}，逼近布雷顿森林鼎盛期的约 3.8 万吨；2025 年央行净购金 \textbf{850 吨}；金价 2026 年 1 月一度接近 \textbf{5600 美元/盎司}。2025 年单一最大买家不是哪个国家，是稳定币发行商 Tether，全年买了超过 100 吨。

\end{databox}

看到这组数字，很多人的第一反应是要回金本位了。这个推论把储备和发钞基础混成了一件事。

央行买黄金，买的是一份保险：黄金不是任何国家的负债，冻结不了，也稀释不掉。但没有一家央行承诺你能拿本币换它。储备是资产端的选择，发钞靠的是负债端的信用。

\section{谁欠你的}

下面这一段是全书的伏笔，读慢一点。

\textbf{每一种钱都是某个人的负债。}你说自己"有钱"，准确的说法是"有人欠你钱"。剩下的问题只有一个：欠你的是谁。

\begin{tablewrap}
\begin{xltabular}{\linewidth}{>{\raggedright\arraybackslash}p{0.20\linewidth}XXX}
\toprule
\bthead{你手里的东西} & \bthead{它本质上是} & \bthead{谁欠你} & \bthead{出事了谁兜} \\
\midrule
\textbf{口袋里的现金} & 中央银行的负债 & 中国人民银行 & 境内最终的钱，不用别人兜 \\
\textbf{银行卡里的存款} & 商业银行的负债 & 你的开户行 & 存款保险，50 万元以内本息全额 \\
\textbf{微信 / 支付宝余额} & 支付机构的负债 & 支付机构 & 客户备付金 100\% 集中交存央行 \\
\textbf{数字人民币} & 中央银行的\textbf{直接}负债 & 中国人民银行 & 等同于现金 \\
\textbf{货币基金份额} & 基金资产的份额，不是存款 & 严格说没人"欠"你 & 没有存款保险，净值理论上会波动 \\
\textbf{美元稳定币} & 发行公司的负债 & 那家发行公司 & 看它的储备资产和监管框架 \\
\bottomrule
\end{xltabular}
\end{tablewrap}

老陈的手机里有三笔钱。银行卡里 27 万，是那家银行欠他的。微信零钱 4300 块，是支付机构欠他的，对应的备付金 100\% 交存在央行。抽屉里压着 1200 块现金，是央行欠他的。

三笔钱在他眼里长得一模一样，法律上隔着三层。

这也解释了数字人民币为什么不是"微信支付的国家队版本"。微信余额是支付机构欠你的，数字人民币是央行直接欠你的，中间少了一层机构信用。

整个体系是一座用负债搭起来的塔。塔尖是央行的负债（现金和准备金），中间是商业银行的负债（存款），再往下是各类机构的负债（余额、稳定币）。越往上越安全，越往下越好用。

老陈抽屉里那张一百块，写在央行的负债栏里。

\begin{chainflow}
\chainnode{央行的负债}\chainarrow \chainnode{你的现金}\chainarrow \chainnode{银行的负债}\chainarrow \chainnode{你的存款}\chainarrow \chainnode{机构的负债}\chainarrow \chainnode{你的余额}
\end{chainflow}

\section{比特币卡在哪一步}

钱只是记账，那能不能不要国家那本账？\term{比特币}{Bitcoin}\index{004@比特币}要回答的就是这个：用密码学和一个分布式网络，让一群互不信任的人维护同一本账。

技术上它做到了。

当日常货币它不行，卡在三个地方。

\begin{flowlist}
  \flowitem{01}{波动}{一天涨跌 5\% 是常态。老陈的一碗牛肉面标 0.00021 个比特币，明天早上他得重新贴一次价签，后天再贴一次。}
  \flowitem{02}{结算}{链上要等确认，手续费还随网络拥堵跳。买碗面付几美元手续费，这笔账算不过来。}
  \flowitem{03}{记账单位失效，最致命的一条}{就算在币圈内部，人们也是拿美元给比特币报价，不是拿比特币给面包报价。它没坐上度量衡那把椅子。另外 2100 万枚的上限意味着，它要真成了主流货币，经济在长它不长，结果是持续通缩，金本位那个老毛病的加强版。}
\end{flowlist}

这里讨论的是它能不能当日常货币，不是它值多少钱。本书不预测任何资产涨跌，也不给买卖建议。

\term{稳定币}{stablecoin}\index{065@稳定币}走的是反方向。它放弃自己当货币，改成锚美元，1 枚兑 1 美元。发行方收你的美元，拿去买短期美债和回购，在链上给你发一枚代币。

拿上面那张表看它：\hlmark{稳定币就是某个私人机构的负债}。你把美元借给一家公司，换回一张能在链上转来转去的欠条。技术是新的，法律结构老得像 19 世纪的私人银行券。

老陈的儿子做跨境电商，收过 USDT。他不关心这些，他只关心到账是三秒还是三天。

\begin{databox}{2026 数据}
美国《GENIUS Act》2025 年 7 月签署生效，给美元稳定币建了联邦层面的监管框架，核心是储备资产的质量和透明度。到 2026 年 8 月，全球稳定币市值约 \textbf{3000 亿美元}量级。

这数字听着大，放进坐标系看：中国广义货币 M2 是三百多万亿元量级。稳定币真正的战场不是取代存款，是跨境支付和 7×24 结算。传统跨境汇款要两三天，扒好几层手续费，链上几秒钟就到。

\end{databox}

\begin{warnbox}{小心}
稳定币稳的是价格锚，不是风险。2023 年 3 月硅谷银行倒掉的时候，某主流美元稳定币因为有一部分储备存在那家银行，价格一度掉到 0.87 美元。\textbf{链上的欠条，安全性取决于链下的资产。}

\end{warnbox}

\bookbreak

回到老陈那碗面。

他敢把面端上来，是因为他相信那行数字明天还有人认。他之所以相信，是因为国家拿它收税、法院认它清偿、整条街都在用它，还有一家央行盯着它的购买力。

雅浦岛现在还留着几千块 Rai。绝大多数从被凿出来那天起就没挪过窝，连沉在海底的那块也还在账上。

\begin{takeaway}{本章一句话}
钱不是财富本身，是一套记录"谁欠谁多少"的公共账本。它有价值不是因为背后有黄金，是因为税收、法律、网络和信誉在同时撑着。

\begin{itemize}
  \item 先有赊账和记账，后有铸币。物物交换起源说是功能论证，不是历史。
  \item 三个功能里记账单位最根本。商家不再拿它标价，货币才算真死。
  \item 每一种钱都是某人的负债。现金是央行的，存款是银行的，余额是支付机构的，稳定币是发行公司的。
\end{itemize}

\end{takeaway}

\begin{bridgebox}
\textbf{下一章}既然你卡里的钱只是银行欠你的，那这笔钱现在人在哪儿？打开银行金库，为什么找不到属于你的那一沓？
\end{bridgebox}

